\chapter{Boundary Gradient Estimates}

Let
\begin{equation}
Qu=a^{ij}(x,u,Du)D_{ij}u+b(x,u,Du)=0
\tag{14.1}
\end{equation}
be elliptic in a domain $\Omega\subset\mathbb R^n$. Write $\lambda,\Lambda$ for
the minimum and maximum eigenvalues of $[a^{ij}]$, and set
\[
\mathcal E=a^{ij}p_ip_j,\qquad
\mathcal E^*=\frac{\mathcal E}{|p|^2},\qquad
\mathcal T=\operatorname{trace}[a^{ij}].
\]

\begin{definition}[Boundary barriers]
\label{gt14-def-boundary-barriers}
\dcref{ch14:14.2}
\group{ch14-boundary-gradient}
For the comparison argument, consider
\begin{equation}
Qu=a^{ij}(x,Du)D_{ij}u+b(x,u,Du),
\tag{14.2}
\end{equation}
where $b$ is nonincreasing in its scalar argument. Let
$u\in C^2(\Omega)\cap C^0(\overline\Omega)$ solve $Qu=0$. At
$x_0\in\partial\Omega$, upper and lower barriers are functions
$w^\pm\in C^2(\mathcal N\cap\Omega)\cap C^1(\mathcal N\cap\overline\Omega)$
such that
\[
\pm Qw^\pm<0\text{ in }\mathcal N\cap\Omega,\qquad
w^\pm(x_0)=u(x_0),\qquad
w^-\leq u\leq w^+\text{ on }\partial(\mathcal N\cap\Omega).
\]
\end{definition}

\begin{proposition}[Barrier derivative comparison]
\label{gt14-prop-barrier-derivative-comparison}
\dcref{ch14:14.3}
\group{ch14-boundary-gradient}
Under the hypotheses of the preceding definition,
$w^-\leq u\leq w^+$ throughout $\mathcal N\cap\Omega$. When the indicated
normal derivatives exist,
\begin{equation}
\frac{\partial w^-}{\partial\nu}(x_0)
\leq\frac{\partial u}{\partial\nu}(x_0)
\leq\frac{\partial w^+}{\partial\nu}(x_0).
\tag{14.3}
\end{equation}
Uniformly bounded barrier gradients therefore give a boundary gradient bound.
\end{proposition}
\begin{proof}
\sketch Apply the comparison principle to $w^-$, $u$, and $w^+$ in
$\mathcal N\cap\Omega$. Divide the resulting inequalities along the inward
normal ray by the distance from $x_0$ and pass to the limit.
\end{proof}

Two transformations will be used repeatedly. If $u=\psi(v)$ with
$\psi\in C^2(I)$ and $\psi'\neq0$, then
\begin{equation}
Qu=\psi'a^{ij}D_{ij}v+\frac{\psi''}{(\psi')^2}\mathcal E+b,
\tag{14.4}
\end{equation}
where the coefficient arguments are $(x,\psi(v),\psi'Dv)$. If
$u=v+\varphi$, then
\begin{equation}
\widetilde Qv=Qu=a^{ij}D_{ij}v+a^{ij}D_{ij}\varphi+b.
\tag{14.5}
\end{equation}
Define
\begin{equation}
\mathcal F(x,z,p,q)=a^{ij}(x,z,p)(p_i-q_i)(p_j-q_j).
\tag{14.6}
\end{equation}
For the transformed operator,
\begin{equation}
\widetilde{\mathcal E}(x,v,Dv)=\mathcal F(x,u,Du,D\varphi).
\tag{14.7}
\end{equation}

\section*{14.1. General Domains}

Suppose an exterior ball $B_R(y)$ is tangent to $\partial\Omega$ at $x_0$, put
$d(x)=\dist(x,\partial B_R(y))$, and take $w=\psi(d)$ with $\psi'>0$. Then
\begin{equation}
\overline Qw
\leq\frac{n-1}{R}\psi'\Lambda+b
+\frac{\psi''}{(\psi')^2}\mathcal E.
\tag{14.8}
\end{equation}
Assume that a nondecreasing $\mu$ satisfies
\begin{equation}
|p|\Lambda+|b|\leq\mu(|z|)\mathcal E
\qquad\text{when }|p|\geq\mu(|z|).
\tag{14.9}
\end{equation}
Writing $M=\sup_\Omega|u|$ and
$\nu=(1+(n-1)/R)\mu(M)$, (14.8) gives
\begin{equation}
\overline Qw\leq
\left(\frac{\psi''}{(\psi')^2}+\nu\right)\mathcal E
\tag{14.10}
\end{equation}
provided $\psi'\geq\mu(M)$. The choice
\begin{equation}
\psi(d)=\frac1\nu\log(1+kd)
\tag{14.11}
\end{equation}
satisfies $\psi''=-\nu(\psi')^2$. For a collar of width $a$, require
\begin{equation}
\psi(a)=M,\qquad ka=e^{\nu M}-1,
\tag{14.12}
\end{equation}
and observe
\begin{equation}
\psi'(d)=\frac{k}{\nu(1+kd)}
\geq\frac{k}{\nu e^{\nu M}}\geq\mu(M)
\quad\text{if }k\geq\mu(M)\nu e^{\nu M}.
\tag{14.13}
\end{equation}
Thus it suffices that
\begin{equation}
k\geq\mu(M)\nu e^{\nu M},\qquad ka=e^{\nu M}-1.
\tag{14.14}
\end{equation}
For zero boundary values the resulting barriers yield
\begin{equation}
|Du(x_0)|\leq\psi'(0)=\mu(M)e^{\nu M}.
\tag{14.15}
\end{equation}

For $u=\varphi$ on $\partial\Omega$, the transformed structure condition is
\begin{equation}
\bigl(|p-D\varphi|+|D^2\varphi|\bigr)\Lambda+|b|
\leq\bar\mu(|z|)\mathcal F(x,z,p,D\varphi)
\tag{14.16}
\end{equation}
when $|p-D\varphi|\geq\bar\mu(|z|)$. The estimate
\begin{equation}
\mathcal F(x,z,p,q)
\geq\frac12\mathcal E-a^{ij}q_iq_j
\geq\frac12\mathcal E-\Lambda|q|^2
\tag{14.17}
\end{equation}
shows that (14.9) implies (14.16) after enlarging $\bar\mu$ in terms of
$|\varphi|_2$.

\begin{theorem}[Boundary gradient bound under an exterior sphere condition]
\label{gt14-boundary-gradient-14-1}\dcref{ch14:14.1}\group{ch14-boundary-gradient}
\source{gilbarg-trudinger:page-0349}
Let $u\in C^2(\Omega)\cap C^1(\overline\Omega)$ satisfy $Qu=0$ in $\Omega$
and $u=\varphi$ on $\partial\Omega$. Suppose $\Omega$ satisfies a uniform
exterior sphere condition, $\varphi\in C^2(\overline\Omega)$, and (14.9) holds.
Then
\begin{equation}
|Du|\leq C\qquad\text{on }\partial\Omega,
\tag{14.18}
\end{equation}
where $C=C(n,M,\mu(M),\Phi,\delta)$,
$M=\sup_\Omega|u|$, $\Phi=|\varphi|_{2;\Omega}$, and $\delta$ is the uniform
exterior-sphere radius.
\end{theorem}
\begin{proof}
\sketch Apply the logarithmic barriers (14.11)--(14.14) to $u-\varphi$.
The transformation (14.5), estimate (14.17), and (14.9) verify (14.16).
The normal derivative comparison (14.3), followed by the tangential derivative
bound from the boundary data, gives (14.18).
\end{proof}

Condition (14.9) may equivalently be expressed as
\begin{equation}
|p|\Lambda,\ |b|=O(\mathcal E)\qquad(|p|\to\infty),
\tag{14.19}
\end{equation}
uniformly for $(x,z)\in\Omega\times(-N,N)$. In particular, it holds when
$\Lambda=O(\lambda)$ and $b=O(\lambda|p|^2)$ on such slabs.

\section*{14.2. Convex Domains}

At a boundary point with an exterior supporting plane, take $d$ to be the
distance from that plane. Then
\begin{equation}
\overline Qw=b+\frac{\psi''}{(\psi')^2}\mathcal E.
\tag{14.20}
\end{equation}
The zero-boundary barrier works under
\begin{equation}
|b|\leq\mu(|z|)\mathcal E
\qquad\text{for }|p|\geq\mu(|z|),
\tag{14.21}
\end{equation}
and nonzero data require
\begin{equation}
\Lambda|D^2\varphi|+|b|
\leq\widetilde\mu(|z|)\mathcal F
\qquad\text{for }|p-D\varphi|\geq\widetilde\mu(|z|).
\tag{14.22}
\end{equation}

\begin{theorem}[Boundary gradient bound in convex domains]
\label{gt14-boundary-gradient-14-2}
\dcref{ch14:14.2}
\group{ch14-boundary-gradient}
\source{gilbarg-trudinger:page-0350}
Let $u,\varphi\in C^2(\Omega)\cap C^1(\overline\Omega)$ satisfy $Qu=0$ and
$u=\varphi$ on $\partial\Omega$. If $\Omega$ is convex and (14.22) holds, then
\begin{equation}
|Du|\leq C\qquad\text{on }\partial\Omega,
\tag{14.23}
\end{equation}
where
$C=C(n,M,\widetilde\mu(M),|\varphi|_{1;\Omega})$ and
$M=\sup_\Omega|u|$.
\end{theorem}
\begin{proof}
\sketch Use the supporting-plane formula (14.20), the logarithmic profile
(14.11), and the transformed condition (14.22). Uniformity follows from
convexity and the displayed dependencies.
\end{proof}

Either
\begin{equation}
\Lambda=o(\mathcal E),\qquad b=O(\mathcal E)
\qquad(|p|\to\infty),
\tag{14.24}
\end{equation}
or
\begin{equation}
\Lambda,\ b=O(\lambda|p|^2)
\qquad(|p|\to\infty)
\tag{14.25}
\end{equation}
implies (14.22), using
\begin{equation}
\mathcal F(x,z,p,q)\geq\lambda(x,z,p)|p-q|^2.
\tag{14.26}
\end{equation}

\begin{corollary}[Convex-domain structural alternatives]
\label{gt14-boundary-gradient-14-3}\dcref{ch14:14.3}\group{ch14-boundary-gradient}
\source{gilbarg-trudinger:page-0350}
Under the solution and boundary hypotheses of Theorem 14.2, suppose
$\varphi\in C^2(\overline\Omega)$ and either (14.24) or (14.25) holds. Then
\begin{equation}
|Du|\leq C\qquad\text{on }\partial\Omega,
\tag{14.27}
\end{equation}
where $C=C(n,M,\widetilde\mu,|\varphi|_{2;\Omega})$.
\end{corollary}
\begin{proof}
\sketch Equations (14.17) and (14.26) reduce the two alternatives to (14.22);
apply Theorem 14.2.
\end{proof}

For the minimal-surface operator
\begin{equation}
\mathfrak Mu=(1+|Du|^2)\Delta u-D_iuD_juD_{ij}u,
\tag{14.28}
\end{equation}
$\lambda=1$ and $\Lambda=1+|p|^2$, so Corollary 14.3 applies on convex
domains.

Suppose now that $\Omega$ lies in a ball $B_R(y)$ tangent at the boundary point.
For the distance $d$ from $\partial B_R(y)$,
\begin{equation}
\overline Qw
\leq-\frac{\psi'}R(\mathcal T-\mathcal E^*)+b
+\frac{\psi''}{(\psi')^2}\mathcal E.
\tag{14.29}
\end{equation}
Assume
\begin{equation}
|a^{ij}D_{ij}\varphi+b|
\leq\frac1R|p-D\varphi|\mathcal T+\bar\mu\mathcal E
\quad\text{for }|p-D\varphi|\geq\bar\mu.
\tag{14.30}
\end{equation}

\begin{theorem}[Uniformly convex-domain estimate]
\label{gt14-boundary-gradient-14-4}\dcref{ch14:14.4}\group{ch14-boundary-gradient}
\source{gilbarg-trudinger:page-0351}
Let $u,\varphi\in C^2(\Omega)\cap C^1(\overline\Omega)$ satisfy $Qu=0$ and
$u=\varphi$ on $\partial\Omega$. If $\Omega$ is uniformly convex with enclosing
sphere radius $R$ and (14.30) holds, then
\begin{equation}
|Du|\leq C\qquad\text{on }\partial\Omega,
\tag{14.31}
\end{equation}
where $C=C(n,M,\bar\mu(M),R,|\varphi|_{1;\Omega})$.
\end{theorem}
\begin{proof}
\sketch Insert the logarithmic profile into (14.29). The negative tangential
trace term and (14.30) make the resulting functions upper and lower barriers;
then use (14.3).
\end{proof}

Condition (14.30) follows from either
\begin{equation}
b=o(\lambda|p|)+O(\lambda|p|^2)
\qquad(|p|\to\infty),
\tag{14.32}
\end{equation}
or
\begin{equation}
\Lambda=O(\lambda|p|^2),\qquad
|b|\leq\frac{|p|}{R}\mathcal T+O(\lambda|p|^2).
\tag{14.33}
\end{equation}

\begin{corollary}[Uniformly convex structural alternatives]
\label{gt14-boundary-gradient-14-5}\dcref{ch14:14.5}\group{ch14-boundary-gradient}
\source{gilbarg-trudinger:page-0352}
Under the solution hypotheses of Theorem 14.4, assume
$\varphi\in C^2(\overline\Omega)$ and either (14.32) or (14.33). Then
\begin{equation}
|Du|\leq C\qquad\text{on }\partial\Omega,
\tag{14.34}
\end{equation}
where $C=C(n,M,\bar\mu(M),|\varphi|_2,R)$.
\end{corollary}
\begin{proof}
\sketch Each displayed asymptotic condition implies (14.30); apply Theorem 14.4.
\end{proof}

For the prescribed mean-curvature equation
\begin{equation}
\mathfrak Mu=nH(x,u,Du)(1+|Du|^2)^{3/2},
\tag{14.35}
\end{equation}
the uniformly convex estimate holds if
\begin{equation}
|H|\leq\frac{n-1}{nR}
\qquad\text{for }|p|\geq\mu(|z|).
\tag{14.36}
\end{equation}

\begin{remark}[Local boundary geometry]
\label{gt14-boundary-gradient-14-r1}\dcref{ch14:14.1}\group{ch14-boundary-gradient}
\source{gilbarg-trudinger:page-0352}
Let $\kappa(x_0)$ be the least principal curvature and $\nu(x_0)$ the inner
normal at $x_0\in\partial\Omega$. A boundary gradient estimate follows if, for
each $x_0$, some neighborhood and nondecreasing $\bar\mu$ satisfy
\begin{equation}
|a^{ij}D_{ij}\varphi+b|
<\kappa(x_0)\mathcal T|p|+\bar\mu(|z|)\mathcal F
\tag{14.37}
\end{equation}
whenever $p-D\varphi$ is sufficiently large and its direction is sufficiently
close to $\pm\nu(x_0)$. Then
\begin{equation}
|Du|\leq C\qquad\text{on }\partial\Omega.
\tag{14.38}
\end{equation}
If $\kappa\geq\kappa_0$ globally, one may use the nonstrict inequality with
$\kappa_0$. For example, for
\begin{equation}
Qu=D_{11}u+(1+|D_2u|^N)D_{22}u=0,\qquad N\geq0,
\tag{14.39}
\end{equation}
a boundary gradient estimate holds for arbitrary $C^2$ boundary data on a
convex planar domain when the curvature is positive except possibly where the
tangent is parallel to the $x_1$-axis.
\end{remark}

The conclusions of this section depend only on the structural conditions in a
neighborhood of $\partial\Omega$.

\section*{14.3. Boundary Curvature Conditions}

Assume $\partial\Omega\in C^2$, let $d=\dist(\cdot,\partial\Omega)$ in a smooth
boundary collar, and put $w=\psi(d)$. Then
\begin{equation}
\overline Qw
=\psi'a^{ij}D_{ij}d+b+\frac{\psi''}{(\psi')^2}\mathcal E.
\tag{14.40}
\end{equation}
For the minimal-surface coefficients
\begin{equation}
a^{ij}=(1+|p|^2)\delta_{ij}-p_ip_j,
\tag{14.41}
\end{equation}
one has
\begin{equation}
a^{ij}D_{ij}d=(1+\psi'^2)\Delta d
\leq-(n-1)(1+\psi'^2)H',
\tag{14.42}
\end{equation}
where $H'$ is the boundary mean curvature at the nearest point.

For the general result, assume the homogeneous decomposition
\begin{equation}
a^{ij}=\lambda a_\infty^{ij}+a_0^{ij},
\qquad
b=|p|\lambda b_\infty+b_0,
\tag{14.43}
\end{equation}
where $a_\infty^{ij}$ and $b_\infty$ are homogeneous of degree zero in $p$,
$[a_\infty^{ij}]$ is nonnegative, and $b_\infty$ is nonincreasing in $z$.
At $y\in\partial\Omega$, with inner normal $\nu$, principal curvatures
$\kappa_1,\ldots,\kappa_{n-1}$, and diagonal entries $a_i(y,\pm\nu)$ in
principal coordinates, define
\begin{equation}
\mathcal K^\pm(y)=\sum_{i=1}^{n-1}a_i(y,\pm\nu)\kappa_i.
\tag{14.44}
\end{equation}
Equivalently,
\begin{equation}
\mathcal K^\pm(y)
=-a_\infty^{ij}(y,\pm Dd(y))D_{ij}d(y).
\tag{14.45}
\end{equation}

The upper-barrier curvature condition is
\begin{equation}
\mathcal K^+(y)\geq b_\infty(y,u(y),\nu(y)),
\tag{14.46}
\end{equation}
and the lower-barrier condition is
\[
\mathcal K^-(y)\geq-b_\infty(y,u(y),-\nu(y)).
\tag{14.46'}
\]
Assume $a_\infty^{ij},b_\infty$ are $C^1$ in a boundary collar and
\begin{equation}
\Lambda,\ |p|\,|a_0^{ij}|,\ |b_0|=O(\mathcal E)
\qquad(|p|\to\infty).
\tag{14.47}
\end{equation}
Equivalently, for a nondecreasing $\mu$,
\begin{equation}
\Lambda+|p|\sum_{i,j}|a_0^{ij}|+|b_0|
\leq\mu(|z|)\mathcal E
\quad\text{when }|p|\geq\mu(|z|).
\tag{14.48}
\end{equation}
The coefficient-modulus constant used below is
\begin{equation}
K=\sup_{x}
\frac{\left|
(a_\infty^{ij}(x,\nu)-a_\infty^{ij}(y,\nu))D_{ij}d(y)
+b_\infty(x,\varphi(y),\nu)-b_\infty(y,\varphi(y),\nu)
\right|}{|x-y|},
\tag{14.49}
\end{equation}
where $y$ is the nearest boundary point to $x$. For nonzero boundary data,
assume
\begin{equation}
\Lambda+|p|\sum_{i,j}|a_0^{ij}|+|b_0|
\leq\bar\mu(|z|)\mathcal F(x,z,p,D\varphi)
\quad\text{when }|p-D\varphi|\geq\bar\mu(|z|).
\tag{14.50}
\end{equation}

\begin{theorem}[Boundary curvature estimate]
\label{gt14-boundary-gradient-14-6}\dcref{ch14:14.6}\group{ch14-boundary-gradient}
\source{gilbarg-trudinger:page-0355}
Let $u\in C^2(\Omega)\cap C^1(\overline\Omega)$ and
$\varphi\in C^2(\overline\Omega)$ satisfy $Qu=0$ and
$u=\varphi$ on $\partial\Omega$. Suppose (14.43), (14.50), and
\begin{equation}
\mathcal K^+\geq b_\infty(y,\varphi(y),\nu),\qquad
\mathcal K^-\geq-b_\infty(y,\varphi(y),-\nu)
\tag{14.51}
\end{equation}
hold at each boundary point. Then
\begin{equation}
|Du|\leq C\qquad\text{on }\partial\Omega,
\tag{14.52}
\end{equation}
where
$C=C(n,M,\bar\mu(M),\Omega,K,|\varphi|_{2;\Omega})$,
$M=\sup_\Omega|u|$, and $K$ is (14.49).
\end{theorem}
\begin{proof}
\sketch Lemmas 14.16--14.17 and (14.45)--(14.51) bound the leading
curvature expression by $Kd$. The logarithmic profile (14.11), with a sufficiently
thin collar, makes $\overline Qw\leq0$ for the upper barrier and gives the lower
barrier after reversing the normal. Apply (14.3) to $u-\varphi$.
\end{proof}

Condition (14.50) follows from either
\begin{equation}
\Lambda=o(\mathcal F),\qquad
|p|a_0^{ij},\ b_0=O(\mathcal F),
\tag{14.53}
\end{equation}
or
\begin{equation}
\Lambda,\ |p|a_0^{ij},\ b_0=O(\lambda|p|^2).
\tag{14.54}
\end{equation}

\begin{corollary}[Boundary curvature structural alternatives]
\label{gt14-boundary-gradient-14-7}\dcref{ch14:14.7}\group{ch14-boundary-gradient}
\source{gilbarg-trudinger:page-0357}
Under the solution hypotheses of Theorem 14.6, assume (14.43), (14.51), and
either (14.53) or (14.54). Then
\begin{equation}
|Du|\leq C\qquad\text{on }\partial\Omega,
\tag{14.55}
\end{equation}
where $C=C(n,M,\bar\mu,\Omega,K,|\varphi|_{2;\Omega})$.
\end{corollary}
\begin{proof}
\sketch Reduce the structural alternatives to (14.50), then apply Theorem 14.6.
\end{proof}

For
\begin{equation}
\mathfrak Mu=nH(x,u)(1+|Du|^2)^{3/2},
\tag{14.56}
\end{equation}
with $H\in C^1(\overline\Omega\times\mathbb R)$ and $D_zH\geq0$, the curvature
conditions become the following result.

\begin{corollary}[Prescribed mean-curvature boundary estimate]
\label{gt14-boundary-gradient-14-8}\dcref{ch14:14.8}\group{ch14-boundary-gradient}
\source{gilbarg-trudinger:page-0357}
Let $u\in C^2(\Omega)\cap C^1(\overline\Omega)$ solve (14.56) with
$u=\varphi\in C^2(\overline\Omega)$ on $\partial\Omega$. If
\begin{equation}
H'(y)\geq\frac n{n-1}|H(y,\varphi(y))|
\qquad(y\in\partial\Omega),
\tag{14.57}
\end{equation}
then
\begin{equation}
|Du|\leq C\qquad\text{on }\partial\Omega,
\tag{14.58}
\end{equation}
where
$C=C(n,M,\Omega,H_1,|\varphi|_{2;\Omega})$ and
$H_1=\sup_{\Omega\times(-M,M)}(|H|+|DH|)$.
\end{corollary}
\begin{proof}
\sketch For the minimal-surface principal matrix,
$\mathcal K^\pm=(n-1)H'$, while
$b_\infty=\mp nH$ in the two normal directions. Thus (14.57) is (14.51);
apply Corollary 14.7.
\end{proof}

The estimate without control of $\Lambda$ uses the stricter assumptions
\begin{equation}
a_0^{ij}=o(\Lambda),\qquad b_0=o(|p|\Lambda)
\qquad(|p|\to\infty)
\tag{14.59}
\end{equation}
and
\begin{equation}
\mathcal K^+>b_\infty(y,\varphi(y),\nu),\qquad
\mathcal K^->-b_\infty(y,\varphi(y),-\nu).
\tag{14.60}
\end{equation}

\begin{theorem}[Strict boundary curvature estimate]
\label{gt14-boundary-gradient-14-9}\dcref{ch14:14.9}\group{ch14-boundary-gradient}
\source{gilbarg-trudinger:page-0358}
Let $u\in C^2(\Omega)\cap C^1(\overline\Omega)$ and
$\varphi\in C^2(\overline\Omega)$ satisfy $Qu=0$ and
$u=\varphi$ on $\partial\Omega$. If (14.43), (14.59), and (14.60) hold at
every boundary point, then
\begin{equation}
|Du|\leq C\qquad\text{on }\partial\Omega,
\tag{14.61}
\end{equation}
where
$C=C(n,M,a_\infty^{ij},a_0^{ij},b_\infty,b_0,\Omega,
|\varphi|_{2;\Omega})$.
\end{theorem}
\begin{proof}
\sketch In a fixed boundary collar, strictness in (14.60), continuity, and
Lemma 14.17 give a uniform negative margin. The linear barriers $w^\pm=\pm kd$
therefore satisfy the strict comparison inequalities for large $k$ by (14.59).
Choose $k$ also to dominate the oscillation of $u-\varphi$, then apply (14.3).
\end{proof}

\section*{14.4. Non-Existence Results}

\begin{theorem}[Comparison with an infinite normal derivative]
\label{gt14-boundary-gradient-14-10}\dcref{ch14:14.10}\group{ch14-boundary-gradient}
\source{gilbarg-trudinger:page-0359}
Let $\Omega$ be bounded and $\Gamma$ a relatively open $C^1$ portion of
$\partial\Omega$. Let $Q$ be elliptic of the form (14.2). If
$u\in C^0(\overline\Omega)\cap C^2(\Omega\cup\Gamma)$ and
$v\in C^0(\overline\Omega)\cap C^2(\Omega)$ satisfy
$Qu\geq Qv$ in $\Omega$, $u\leq v$ on
$\partial\Omega\setminus\Gamma$, and
$\partial v/\partial\nu=-\infty$ on $\Gamma$, then $u\leq v$ in $\Omega$.
\end{theorem}
\begin{proof}
\sketch Theorem 10.1 bounds $\sup_\Omega(u-v)$ by its positive boundary
maximum on $\Gamma$. Since
$\partial_\nu(u-v)=+\infty$ there, such a maximum cannot occur; hence
$u-v\leq0$.
\end{proof}

For $y\in\partial\Omega$, $\delta=\diam\Omega$, and
$w=m+\psi(r)$, $r=|x-y|$, with
$\psi(\delta)=0$, $\psi'\leq0$, and $\psi'(a)=-\infty$,
\begin{equation}
\overline Qw
=\frac{\psi'}r(\mathcal T-\mathcal E^*)
+\frac{\psi''}{(\psi')^2}\mathcal E+b.
\tag{14.62}
\end{equation}
Whenever this is negative on $\{r>a,\ |u|>M\}$, Theorem 14.10 gives
\begin{equation}
\sup_{\Omega\setminus B_a(y)}u\leq M+m+\psi(a),
\qquad
m=\sup_{\partial\Omega\setminus B_a(y)}u^+.
\tag{14.63}
\end{equation}

One sufficient condition is
\begin{equation}
b\leq-|p|^\theta\mathcal E
\tag{14.64}
\end{equation}
for $|z|\geq M$, $|p|\geq L$, with profile
\begin{equation}
\psi(r)=K\bigl((\delta-a)^\beta-(r-a)^\beta\bigr),
\qquad \beta=\frac{\theta}{1+\theta}.
\tag{14.65}
\end{equation}
Another is
\begin{equation}
b\leq0,\qquad
\mathcal E\leq\mu(\mathcal T-\mathcal E^*)|p|^{1-\theta},
\tag{14.66}
\end{equation}
with
\begin{equation}
\psi(r)=\left(\frac\mu\beta\right)^\beta
\int_r^\delta\left(\log\frac ta\right)^{-\beta}\,dt,
\qquad \beta=\frac1{1+\theta},
\tag{14.67}
\end{equation}
for which $\psi(a)\to0$ as $a\downarrow0$.

If an interior ball $B_R(x_0)$ is tangent at $y$, use
\begin{equation}
w^*(x)=m^*+\chi(r),\qquad r=|x-x_0|,
\tag{14.68}
\end{equation}
and
\begin{equation}
\overline Qw^*
=\frac{\chi'}r(\mathcal T-\mathcal E^*)
+\frac{\chi''}{(\chi')^2}\mathcal E+b.
\tag{14.69}
\end{equation}
Assume, for $0<R'<R$,
\begin{equation}
b+\frac{|p|}{R'}\mathcal T\leq-|p|^\theta\mathcal E^*
\tag{14.70}
\end{equation}
and choose
\begin{equation}
\chi(r)=K\bigl((R-\varepsilon)^\beta
-(R-\varepsilon-r)^\beta\bigr),
\qquad\beta=\frac{\theta}{1+\theta}.
\tag{14.71}
\end{equation}
Then comparison yields
\begin{equation}
\sup_{\widetilde\Omega}u
\leq M+m^*+\chi(R-\varepsilon),
\tag{14.72}
\end{equation}
and, after combining the two barriers,
\begin{equation}
u(y)\leq2M+m+K\bigl((\delta-R)^\beta+R^\beta\bigr).
\tag{14.73}
\end{equation}
Applying the same argument to $-u$ gives the symmetric sufficient condition
\begin{equation}
|b|\geq\frac{|p|}{R'}\mathcal T+|p|^\theta\mathcal E^*.
\tag{14.74}
\end{equation}

\begin{theorem}[Non-solvability from excessive lower-order growth]
\label{gt14-boundary-gradient-14-11}\dcref{ch14:14.11}\group{ch14-boundary-gradient}
\source{gilbarg-trudinger:page-0362}
Let $\Omega$ be bounded and suppose $Q$ satisfies (14.74), with $R$ the radius
of the largest ball contained in $\Omega$. Then some
$\varphi\in C^\infty(\overline\Omega)$ makes the Dirichlet problem
$Qu=0$ in $\Omega$, $u=\varphi$ on $\partial\Omega$, unsolvable.
\end{theorem}
\begin{proof}
\sketch Choose boundary values that violate the universal restriction (14.73).
Any classical solution would satisfy that estimate by Theorem 14.10, a
contradiction.
\end{proof}

For the curvature obstruction, assume (14.43), (14.59), $b_\infty$ is independent
of $z$, and
\begin{equation}
b\leq0,\qquad
\mathcal E\leq\mu\Lambda|p|^{1-\theta}
\tag{14.75}
\end{equation}
for $|z|\geq M$. The maximum principle gives
\begin{equation}
\sup_\Omega u\leq M+\sup_{\partial\Omega}u.
\tag{14.76}
\end{equation}
Suppose at $y\in\partial\Omega$,
\begin{equation}
\mathcal K^-(y)<-b_\infty(y,-\nu(y)).
\tag{14.77}
\end{equation}
There is then an interior tangent quadric $\mathcal S$ whose weighted curvature
satisfies
\begin{equation}
\mathcal K^-_{\mathcal S}(y)\leq\mathcal K^-(y)+\eta.
\tag{14.78}
\end{equation}
For $d=\dist(\cdot,\mathcal S)$, take
\begin{equation}
w^*=m^*+\chi(d)
\tag{14.79}
\end{equation}
with
\begin{equation}
\chi(d)=K\bigl((2a-\varepsilon)^\beta-(d-\varepsilon)^\beta\bigr),
\qquad\beta=\frac{\theta}{1+\theta}.
\tag{14.80}
\end{equation}
Theorem 14.10 gives
\begin{equation}
\sup_{\widetilde\Omega}u
\leq M+m^*+K(2a)^\beta,
\tag{14.81}
\end{equation}
and hence
\begin{equation}
u(y)\leq2M+m+\psi(a)+K(2a)^\beta.
\tag{14.82}
\end{equation}

\begin{theorem}[Non-solvability from a boundary curvature defect]
\label{gt14-boundary-gradient-14-12}
\dcref{ch14:14.12}
\group{ch14-boundary-gradient}
\source{gilbarg-trudinger:page-0364}
Let $\Omega$ be a bounded $C^2$ domain and suppose $Q$ satisfies
(14.43), (14.59), and (14.75). If at some $y\in\partial\Omega$,
\begin{equation}
\mathcal K^-(y)<-b_\infty(y,-\nu(y)),
\tag{14.83}
\end{equation}
then some $\varphi\in C^\infty(\overline\Omega)$ makes the Dirichlet problem
$Qu=0$, $u=\varphi$ on $\partial\Omega$, unsolvable. Replacing $u$ by $-u$
gives the analogous conclusion under
\begin{equation}
\mathcal K^+(y)<b_\infty(y,\nu(y)),
\tag{14.84}
\end{equation}
with $b$ replaced by $-b$ in (14.75). If $M=0$ in (14.76), the boundary
norm of $\varphi$ may be arbitrarily small.
\end{theorem}
\begin{proof}
\sketch The tangent-quadric construction (14.77)--(14.82) imposes a bound on
the value at $y$ in terms of boundary values away from $y$. Choose smooth
boundary data violating that bound. Negating the equation gives the second
alternative.
\end{proof}

\begin{corollary}[Prescribed mean-curvature obstruction]
\label{gt14-boundary-gradient-14-13}\dcref{ch14:14.13}\group{ch14-boundary-gradient}
\source{gilbarg-trudinger:page-0364}
Let $\Omega$ be a bounded $C^2$ domain. Suppose $H\in C^0(\overline\Omega)$
has one sign and, at some $y\in\partial\Omega$,
\begin{equation}
H'(y)<\frac n{n-1}|H(y)|.
\tag{14.85}
\end{equation}
For every $\varepsilon>0$, some
$\varphi\in C^\infty(\overline\Omega)$ with $\sup|\varphi|\leq\varepsilon$
makes the prescribed mean-curvature Dirichlet problem unsolvable.
\end{corollary}
\begin{proof}
\sketch Substitute the prescribed mean-curvature coefficients into Theorem 14.12.
\end{proof}

\begin{theorem}[Jenkins--Serrin criterion]
\label{gt14-boundary-gradient-14-14}\dcref{ch14:14.14}\group{ch14-boundary-gradient}
\source{gilbarg-trudinger:page-0364}
Let $\Omega$ be a bounded $C^{2,\gamma}$ domain, $0<\gamma<1$. The
minimal-surface Dirichlet problem
\[
\mathfrak Mu=0\text{ in }\Omega,\qquad u=\varphi\text{ on }\partial\Omega,
\]
is solvable for every $\varphi\in C^{2,\gamma}(\overline\Omega)$ if and only
if the boundary mean curvature $H'$ is nonnegative everywhere.
\end{theorem}
\begin{proof}
\sketch Sufficiency combines Corollary 14.8 with the existence procedure for
the special case (11.7). If $H'<0$ somewhere, Corollary 14.13 supplies smooth
boundary data for which no solution exists.
\end{proof}

Allowing $b_\infty$ to depend on $z$ replaces (14.85) by
\begin{equation}
H'(y)<\sup_{z\in\mathbb R}\frac n{n-1}|H(y,z)|.
\tag{14.86}
\end{equation}

\section*{14.5. Continuity Estimates}

For continuous boundary data, fix $y\in\partial\Omega$ and choose $\delta>0$ so
that $|\varphi(x)-\varphi(y)|<\varepsilon$ when $|x-y|<\delta$. Define
\begin{equation}
\varphi^\pm(x)=\varphi(y)\pm
\left(\varepsilon+\frac{2\sup|\varphi|}{\delta^2}|x-y|^2\right).
\tag{14.87}
\end{equation}
The previous barriers give
\begin{equation}
\varphi^-+w^-\leq u\leq\varphi^++w^+
\tag{14.88}
\end{equation}
near $y$, and hence
\begin{equation}
|u(x)-\varphi(y)|
\leq\varepsilon+w(x)
+\frac{2\sup|\varphi|}{\delta^2}|x-y|^2.
\tag{14.89}
\end{equation}

\begin{theorem}[Boundary modulus of continuity]
\label{gt14-boundary-gradient-14-15}\dcref{ch14:14.15}\group{ch14-boundary-gradient}
\source{gilbarg-trudinger:page-0365}
Let $u,\varphi\in C^0(\overline\Omega)\cap C^2(\Omega)$ satisfy $Qu=0$ and
$u=\varphi$ on $\partial\Omega$. Suppose $Q$ and $\Omega$ satisfy the structural
and geometric hypotheses of Theorem 14.1, Corollary 14.3, Corollary 14.5,
Corollary 14.7, or Theorem 14.9. Then the boundary modulus of continuity of
$u$ is controlled by that of $\varphi$, $\sup_{\partial\Omega}|\varphi|$,
$\sup_\Omega|u|$, $\Omega$, and the coefficients of $Q$.
\end{theorem}
\begin{proof}
\sketch Use $\varphi^\pm$ from (14.87) as smooth local boundary majorants and
minorants. Applying the relevant barrier theorem gives (14.88)--(14.89);
optimize $\varepsilon$ and $\delta$ using the modulus of $\varphi$.
\end{proof}

If the boundary data are $C^2$ but the solution is only continuous up to the
boundary, the same argument bounds
\[
\sup_{\substack{x\in\Omega\\y\in\partial\Omega}}
\frac{|u(x)-u(y)|}{|x-y|}.
\]

\section*{14.6. Boundary Curvatures and the Distance Function}

Define
\begin{equation}
d(x)=\dist(x,\partial\Omega).
\tag{14.90}
\end{equation}
Then
\begin{equation}
|d(x)-d(y)|\leq|x-y|.
\tag{14.91}
\end{equation}
If $\partial\Omega\in C^2$, choose the inner normal as the $x_n$ direction at
$y_0\in\partial\Omega$ and write the boundary locally as
$x_n=\varphi(x')$. The principal curvatures are the eigenvalues
$\kappa_1,\ldots,\kappa_{n-1}$ of $D^2\varphi(y_0)$, and
\begin{equation}
H(y_0)=\frac1{n-1}\sum_{i=1}^{n-1}\kappa_i
=\frac1{n-1}\Delta\varphi(y_0).
\tag{14.92}
\end{equation}
In principal coordinates,
\begin{equation}
D^2\varphi(y_0)=\diag(\kappa_1,\ldots,\kappa_{n-1}).
\tag{14.93}
\end{equation}
The inner normal along the graph is
\begin{equation}
\nu_i=\frac{-D_i\varphi}{\sqrt{1+|D\varphi|^2}}\ (i<n),
\qquad
\nu_n=\frac1{\sqrt{1+|D\varphi|^2}},
\tag{14.94}
\end{equation}
so
\begin{equation}
D_j\nu_i(y_0)=-\kappa_i\delta_{ij}
\qquad(i,j<n).
\tag{14.95}
\end{equation}

\begin{lemma}[Regularity of the distance function]
\label{gt14-boundary-gradient-14-16}\dcref{ch14:14.16}\group{ch14-boundary-gradient}
\source{gilbarg-trudinger:page-0367}
Let $\Omega$ be bounded and $\partial\Omega\in C^k$, $k\geq2$. There is
$\mu>0$, depending on $\Omega$, such that
$d\in C^k(\Gamma_\mu)$, where
$\Gamma_\mu=\{x\in\Omega:d(x)<\mu\}$.
\end{lemma}
\begin{proof}
\sketch Choose $\mu$ below the uniform interior normal radius. Every
$x\in\Gamma_\mu$ has a unique nearest boundary point $y(x)$ and
\begin{equation}
x=y+\nu(y)d.
\tag{14.96}
\end{equation}
In principal coordinates the normal map $(y',d)\mapsto y+\nu(y)d$ has
\begin{equation}
D\mathbf g=\diag(1-\kappa_1d,\ldots,1-\kappa_{n-1}d,1)
\tag{14.97}
\end{equation}
and
\begin{equation}
\det D\mathbf g=\prod_{i=1}^{n-1}(1-\kappa_id)>0.
\tag{14.98}
\end{equation}
The inverse function theorem makes $y$ and $d$ smooth; since
$Dd=\nu(y(x))$, one gains the final derivative and obtains $d\in C^k$.
\end{proof}

\begin{lemma}[Hessian of the distance function]
\label{gt14-boundary-gradient-14-17}\dcref{ch14:14.17}\group{ch14-boundary-gradient}
\source{gilbarg-trudinger:page-0367}
\source{gilbarg-trudinger:page-0368}
Under the hypotheses of Lemma 14.16, let $x_0\in\Gamma_\mu$ and let
$y_0\in\partial\Omega$ be its nearest boundary point. In principal coordinates
at $y_0$,
\begin{equation}
D^2d(x_0)=\diag\left(
\frac{-\kappa_1}{1-\kappa_1d(x_0)},\ldots,
\frac{-\kappa_{n-1}}{1-\kappa_{n-1}d(x_0)},0\right).
\tag{14.99}
\end{equation}
\end{lemma}
\begin{proof}
\sketch Differentiate $Dd(x)=\nu(y(x))$. Formulae (14.95) and (14.97) give
the tangential diagonal entries; $|Dd|=1$ gives all mixed and normal entries.
\end{proof}

\begin{proposition}[Curvature of a level hypersurface]
\label{gt14-prop-level-hypersurface-curvature}
\dcref{ch14:14.100--14.102}
\group{ch14-boundary-gradient}
Let a $C^2$ hypersurface $\mathcal S$ be locally given by $\psi=0$ with
$|D\psi|>0$, oriented toward positive $\psi$. Its unit normal is
\begin{equation}
\nu=\frac{D\psi}{|D\psi|}.
\tag{14.100}
\end{equation}
In principal coordinates at $y_0$,
\begin{equation}
D_i\left(\frac{D_j\psi}{|D\psi|}\right)
=-\kappa_i\delta_{ij}\quad(i,j<n),
\tag{14.101}
\end{equation}
with the normal-direction derivatives zero. Hence
\begin{equation}
H(y_0)=-\frac1{n-1}
\left.D_i\left(\frac{D_i\psi}{|D\psi|}\right)\right|_{y_0}.
\tag{14.102}
\end{equation}
\end{proposition}
\begin{proof}
\sketch Differentiate the normalized gradient in principal coordinates. Its
tangential eigenvalues are $-\kappa_1,\ldots,-\kappa_{n-1}$ and its normal
eigenvalue is zero; taking the trace gives (14.102).
\end{proof}

\begin{corollary}[Curvatures of a graph]
\label{gt14-cor-graph-curvatures}
\dcref{ch14:14.103--14.105}
\group{ch14-boundary-gradient}
For the graph of $u\in C^2(\Omega)$ in $\mathbb R^{n+1}$,
\begin{equation}
H(x_0)=\frac1n
\left.D_i\left(\frac{D_iu}{\sqrt{1+|Du|^2}}\right)\right|_{x_0}.
\tag{14.103}
\end{equation}
Writing $\nu_i=D_iu/\sqrt{1+|Du|^2}$, the squared norm of the second
fundamental form is
\begin{equation}
\mathcal G^2=\sum_{i=1}^n\kappa_i^2
=\sum_{i,j=1}^nD_i\nu_jD_j\nu_i,
\tag{14.104}
\end{equation}
and the Gauss curvature is
\begin{equation}
K=\prod_{i=1}^n\kappa_i
=\det[D_i\nu_j]
=\frac{\det D^2u}{(1+|Du|^2)^{(n+2)/2}}.
\tag{14.105}
\end{equation}
The graph is minimal precisely when $H=0$.
\end{corollary}
\begin{proof}
\sketch Apply the preceding proposition to the level function
$\psi(x,x_{n+1})=x_{n+1}-u(x)$. Computing the normalized normal derivative,
its trace, and its determinant gives (14.103)--(14.105).
\end{proof}
